\section{Estrategia de análisis de datos}
\label{sec:analisis}

La estrategia evalúa la efectividad del módulo comparando el desempeño de la
señal filtrada frente a la línea base telemétrica degradada, respecto de la
verdad de referencia física.

\subsection{Métricas cuantitativas y coherencia dimensional}

Se cuantifican el error cuadrático medio y su raíz:

\begin{align}
\mathrm{MSE}_{\text{SCADA}} &= \frac{1}{N}\sum_{k=1}^{N}\left(V_{k}-y_{k}\right)^{2},
&
\mathrm{MSE}_{\text{KF}} &= \frac{1}{N}\sum_{k=1}^{N}
\left(V_{k}-\hat{V}_{k|k}\right)^{2}
\label{eq:mse}\\[0.3em]
\mathrm{RMSE} &= \sqrt{\mathrm{MSE}},
&
\Delta\mathrm{MSE}_{\%} &= \frac{\mathrm{MSE}_{\text{SCADA}}
-\mathrm{MSE}_{\text{KF}}}{\mathrm{MSE}_{\text{SCADA}}}\times 100
\label{eq:reduccion}
\end{align}

Se distingue de manera explícita entre métricas cuadráticas y lineales, punto
observado en la retroalimentación de la Fase 3. Dado que el volumen se expresa
en \si{\metre\cubed}, el error cuadrático medio se expresa necesariamente en
\si{\metre\tothe{6}}, mientras que su raíz recupera las unidades de la variable
física, \si{\metre\cubed}.

En consecuencia, y como criterio de diseño del análisis, \textbf{toda la
inferencia estadística se realiza sobre el RMSE}, de modo que los estadísticos,
los intervalos de confianza y el tamaño del efecto quedan expresados en
\si{\metre\cubed} y son dimensionalmente homogéneos con el estado estimado. El
MSE se reporta como métrica descriptiva y como base del porcentaje de reducción,
que es adimensional. Conviene advertir que ambos porcentajes no son
intercambiables: una reducción del MSE del \SI{89,57}{\percent} corresponde a una
reducción del RMSE del \SI{67,80}{\percent}, por la relación cuadrática entre
ambas métricas. El informe reporta las dos y declara cuál corresponde en cada
caso.

\subsection{Diseño de la prueba de hipótesis}

Cada corrida de Monte Carlo constituye una realización independiente del proceso
estocástico y aporta un par $(\mathrm{RMSE}_{\text{SCADA},m},\
\mathrm{RMSE}_{\text{KF},m})$. Sobre los $M = 50$ pares se define la diferencia
$d_{m} = \mathrm{RMSE}_{\text{SCADA},m} - \mathrm{RMSE}_{\text{KF},m}$ y se
contrasta

\begin{equation}
H_{0}: \mu_{d} = 0 \qquad \text{frente a} \qquad H_{1}: \mu_{d} > 0,
\qquad \alpha = 0{,}05
\label{eq:hipotesis}
\end{equation}

mediante una \textbf{prueba t de Student para muestras pareadas} de una cola. El
pareo es metodológicamente pertinente porque ambas mediciones provienen de la
misma realización del proceso, lo que elimina la variabilidad entre corridas
como fuente de ruido en la comparación. Se reportan además el intervalo de
confianza al \SI{95}{\percent} sobre $\mu_{d}$ y el tamaño del efecto mediante
la $d$ de Cohen para muestras pareadas.

\subsection{Verificación del supuesto de normalidad y contraste de robustez}

La validez de la prueba t descansa sobre el supuesto de normalidad de las
diferencias. Ese supuesto no se asume: se verifica empíricamente mediante la
\textbf{prueba de Shapiro-Wilk} sobre la distribución de $d_{m}$.

Adicionalmente, y con independencia del resultado de esa verificación, se aplica
la \textbf{prueba no paramétrica de Wilcoxon de rangos con signo} como contraste
de robustez. Esta decisión responde a un criterio de rigor: si ambas pruebas
concuerdan, la conclusión no depende del cumplimiento del supuesto paramétrico,
y la inferencia queda blindada frente a la objeción de normalidad. Los
resultados de ambas pruebas se reportan en la sección~\ref{sec:resultados}.

\subsection{Análisis de la variable latente}

Se evalúa por separado el desempeño sobre la tasa de llenado $\dot{V}_{k}$,
variable que ningún sensor mide. La línea base para esta comparación es la
derivada numérica de la señal cruda descrita en la
sección~\ref{sec:procedimientos}, cuyo ruido se amplifica por un factor
$\sqrt{2R}/\Delta t$ respecto del ruido del sensor. Este análisis es el que
permite evaluar directamente la afirmación central sobre observabilidad parcial,
puesto que mide la capacidad del módulo de reconstruir aquello que no se observa.

\subsection{Análisis de sensibilidad y robustez paramétrica}

Se ejecutan dos experimentos complementarios que operacionalizan el criterio de
sensibilidad y robustez declarado en la sección~\ref{sec:calidad}.

El \textbf{barrido de degradación} incrementa la varianza real del ruido de
medición desde condiciones benignas hasta degradación severa, manteniendo el
filtro correctamente sintonizado, y mide cómo escala el aporte del módulo con la
hostilidad del entorno.

La \textbf{malla de desajuste} sintoniza el filtro con covarianzas escaladas por
factores $(\alpha, \beta)$ respecto de las reales, de modo que
$\mathbf{Q}_{\text{filtro}} = \alpha\,\mathbf{Q}_{\text{real}}$ y
$\mathbf{R}_{\text{filtro}} = \beta\,\mathbf{R}_{\text{real}}$. Representa el
escenario realista en que un ingeniero de planta calibra el estimador con
parámetros imprecisos, y cuantifica cuánta degradación de sintonización tolera
el módulo antes de perder utilidad. Todas las combinaciones se evalúan sobre
exactamente los mismos escenarios físicos, de modo que las diferencias
observadas sean atribuibles únicamente a la calibración.

\subsection{Organización de los resultados}

Los resultados se organizan en cuatro bloques de presentación: métricas
descriptivas de error en tabla, comparación pareada e inferencia estadística en
tabla y figura, reconstrucción de series temporales en figura, y análisis de
sensibilidad en figura de barrido y mapa de calor. Esta organización se mantiene
idéntica en el artículo científico, según se documenta en el
Anexo~\ref{anexo:correspondencia}.
